% Part: methods
% Chapter: induction
% Section: inductive-definitions

\documentclass[../../../include/open-logic-section]{subfiles}

\begin{document}

\olfileid{mth}{ind}{idf}

\olsection{विगमनाधारित व्याख्या}

तर्कशास्त्रात अनेकदा वस्तूंचे प्रकार \emph{विगमनाधारित रीतीने} परिभाषित करतो. म्हणजे त्या प्रकारची वस्तू कशाला म्हणायचे याचे नियम देतो; त्यांत त्या प्रकारच्या जुन्या वस्तूंपासून नव्या वस्तू कशा मिळवायच्या हे सांगितलेले असते. उदाहरणार्थ, भाषेतील पदे आणि !!{formula}s यांसारखे चिन्हांचे विशिष्ट अनुक्रम अनेकदा विगमनाने परिभाषित करतो. एक साधे उदाहरण पाहू: $\mathrm{a}$, $\mathrm{b}$, $\mathrm{c}$, $\mathrm{d}$ ही अक्षरे,~$\circ$ हे चिन्ह आणि $[$, $]$ हे कंस यांपासून बनलेल्या चिन्हमाला घ्या. उदाहरणार्थ, “$[[\mathrm{c} \circ \mathrm{d}][$”, “$[\mathrm{a}[]\circ]$”, “$\mathrm{a}$” किंवा “$[[\mathrm{a} \circ \mathrm{b}]\circ \mathrm{d}]$”. पहिल्या दोन चिन्हमालांत काहीतरी “चुकीचे” आहे असे वाटत असेल: पहिल्या चिन्हमालेत कंस मुळीच जुळत नाहीत; आणि “$\circ$” ने स्वतः अर्थपूर्ण असलेल्या अभिव्यक्ती “जोडल्या” पाहिजेत असे वाटेल. तिसरी आणि चौथी चिन्हमाला अधिक नीट दिसतात: प्रत्येक “$[$” साठी बंद करणारा “$]$” आहे---कंस असतील तर---आणि प्रत्येक $\circ$ च्या दोन्ही बाजूंना स्वतः नीट असलेल्या अभिव्यक्ती आहेत, ज्यांभोवती कंसांची जोडी आहे.

“नीटस पद” नेमके कशाला म्हणायचे हे निश्चित करायचे आहे. प्रथम, प्रत्येक स्वतंत्र अक्षर नीटस आहे. एकटे अक्षर नसलेले प्रत्येक नीटस पद “$[t \circ s]$” या रूपाचे असले पाहिजे; येथे $s$ आणि $t$ ही स्वतः नीटस असतात. उलट, $t$ आणि $s$ नीटस असतील, तर त्यांच्यामध्ये $\circ$ ठेवून आणि भोवती कंसांची जोडी घालून नवे नीटस पद बनवता येते. या क्रिया वापरून नीटस पदांचा संच \emph{परिभाषित} करता येईल. ही \emph{विगमनाधारित व्याख्या} आहे.

\begin{defn}[नीटस पदे]
\emph{नीटस पदांचा} संच पुढील विगमनाधारित व्याख्येने दिलेला आहे:
\begin{enumerate}
\item $\mathrm{a}$, $\mathrm{b}$, $\mathrm{c}$, $\mathrm{d}$ यांपैकी कोणतेही अक्षर नीटस पद आहे.
\item $s_1$ आणि $s_2$ नीटस पदे असतील, तर $[s_1 \circ s_2]$ देखील नीटस पद आहे.
\item इतर काहीही नीटस पद नाही.
\end{enumerate}
\end{defn}

या व्याख्येनुसार एखादी गोष्ट नीटस पद असते तेव्हाच आणि फक्त तेव्हा जेव्हा (1) आणि~(2) या दोन अटींनुसार सांत संख्येने पायऱ्यांत ती रचता येते. पहिल्या पायरीत केवळ स्वतंत्र अक्षरे असलेली सर्व नीटस पदे रचतो:
\[
\mathrm{a}, \mathrm{b}, \mathrm{c}, \mathrm{d}
\]
दुसऱ्या पायरीत रचलेल्या पदांना (2) लागू करतो. दोन अक्षरांच्या सर्व जोड्यांसाठी पुढील पदे मिळतील:
\[
[\mathrm{a} \circ \mathrm{a}], [\mathrm{a} \circ \mathrm{b}],
[\mathrm{b} \circ \mathrm{a}], \dots, [\mathrm{d} \circ \mathrm{d}]
\]
तिसऱ्या पायरीत आत्तापर्यंत रचलेल्या कोणत्याही दोन नीटस पदांना पुन्हा (2) लागू करतो. उदाहरणार्थ, $[\mathrm{a} \circ [\mathrm{a} \circ \mathrm{a}]]$ असे नवे नीटस पद मिळते: येथे $t$ हे पायरी~1 मधील $\mathrm{a}$ आहे आणि $s$ हे पायरी~2 मधील $[\mathrm{a} \circ \mathrm{a}]$ आहे. तसेच पायरी~2 मधील $[\mathrm{b} \circ \mathrm{c}]$ आणि $[\mathrm{d} \circ \mathrm{b}]$ या दोन पदांपासून $[[\mathrm{b} \circ \mathrm{c}] \circ [\mathrm{d} \circ \mathrm{b}]]$ रचता येते. असेच पुढे. या रीतीने न रचलेली कोणतीही गोष्ट नीटस पदांच्या संचात येऊ नये, हे खंड (3) निश्चित करतो.

“नीट वाटणारा” प्रत्येक चिन्हानुक्रम या व्याख्येनुसार नीटस आहे, हे अद्याप सिद्ध केलेले नाही, हे लक्षात घ्या. मात्र आपण रचू शकणारी प्रत्येक गोष्ट खरोखर “नीट वाटते”, हे स्पष्ट असले पाहिजे: कंस जुळतात आणि $\circ$ ने स्वतः नीटस असलेले भाग जोडलेले असतात.

विगमनाधारित व्याख्यांचे महत्त्वाचे वैशिष्ट्य असे की सर्व नीटस पदांविषयी काही सिद्ध करायचे असेल, तर कोणती प्रकरणे पाहावी लागतील हे व्याख्या सांगते. उदाहरणार्थ, $t$ हे नीटस पद आहे असे दिले असेल, तर $t$ कसे असू शकते हे व्याख्या सांगते: $t$ हे अक्षर असेल किंवा ते $[s_1 \circ s_2]$ या रूपाचे असेल, जिथे $s_1$ आणि~$s_2$ ही नीटस पदे आहेत. खंड (3) मुळे या दोनच शक्यता आहेत.

विगमनाधारित व्याख्येने दिलेल्या संचातील सर्व वस्तूंविषयी विधाने सिद्ध करताना प्रबल विगमन विशेष महत्त्वाचे ठरते. उदाहरणार्थ,~$n$ लांबीच्या प्रत्येक नीटस पदात $[$ ची संख्या~$< n/2$ आहे हे सिद्ध करायचे आहे असे समजा. याचा सर्व~$n$ विषयीचे विधान म्हणून विचार करता येतो: प्रत्येक $n$ साठी,~$n$ लांबीच्या कोणत्याही नीटस पदात $[$ ची संख्या $< n/2$ आहे.

\begin{prop}
कोणत्याही $n$ साठी,~$n$ लांबीच्या नीटस पदात $[$ ची संख्या $< n/2$ आहे.
\end{prop}

\begin{proof}
प्रबल विगमनाने हा निष्कर्ष सिद्ध करण्यासाठी पुढील सशर्त विधान सत्य आहे हे दाखवावे लागेल:
\begin{quote}
प्रत्येक $l < k$ साठी,~$l$ लांबीच्या कोणत्याही नीटस पदात $[$ ची संख्या $< l/2$ असेल, तर~$k$ लांबीच्या कोणत्याही नीटस पदात $[$ ची संख्या $< k/2$ असते.
\end{quote}
हे सशर्त विधान दाखवण्यासाठी त्याचे पूर्वांग सत्य आहे असे गृहीत धरा. म्हणजे कोणत्याही $l<k$ साठी,~$l$ लांबीच्या नीटस पदात $[$ ची संख्या $< l/2$ आहे असे गृहीत धरा. याला विगमन गृहीतक म्हणतो. ~$k$ लांबीच्या नीटस पदांसाठीही तेच सत्य आहे हे दाखवायचे आहे.

म्हणून $t$ हे~$k$ लांबीचे नीटस पद आहे असे समजा. नीटस पदांची व्याख्या विगमनाधारित असल्यामुळे दोन प्रकरणे आहेत: (1)~$t$ हे स्वतंत्र अक्षर आहे; किंवा (2)~$t$ हे $[s_1 \circ s_2]$ या रूपाचे आहे, जिथे $s_1$ आणि~$s_2$ ही नीटस पदे आहेत.
\begin{enumerate}
\item $t$ हे अक्षर आहे. मग $k = 1$ आणि $t$ मध्ये $[$ ची संख्या~$0$ आहे. $0 < 1/2$ असल्यामुळे विधान खरे आहे.
\item $t$ हे $[s_1 \circ s_2]$ या रूपाचे आहे, जिथे $s_1$ आणि~$s_2$ ही नीटस पदे आहेत. ~$s_1$ ची लांबी $l_1$ आणि~$s_2$ ची लांबी $l_2$ असू द्या. मग $t$ ची लांबी~$k$ ही $l_1+l_2+3$ आहे: $s_1$ आणि $s_2$ यांच्या लांबी अधिक $[$, $\circ$, $]$ ही तीन चिन्हे. $l_1+l_2+3$ हे नेहमी $l_1$ पेक्षा मोठे असल्यामुळे $l_1 < k$. त्याचप्रमाणे $l_2 < k$. म्हणजे $s_1$ आणि~$s_2$ या पदांना विगमन गृहीतक लागू होते: $s_1$ मध्ये $[$ ची संख्या~$m_1$ ही $< l_1/2$ आहे; आणि~$s_2$ मध्ये $[$ ची संख्या~$m_2$ ही $< l_2/2$ आहे.

$t$ मध्ये $[$ ची संख्या म्हणजे~$s_1$ मधील $[$ ची संख्या, अधिक~$s_2$ मधील $[$ ची संख्या, अधिक~$1$; म्हणजे $m_1 + m_2 + 1$. $m_1 < l_1/2$ आणि $m_2 < l_2/2$ असल्यामुळे:
\[
  m_1 + m_2 + 1 < \frac{l_1}{2} + \frac{l_2}{2} + 1 = \frac{l_1+l_2+2}{2} < \frac{l_1+l_2+3}{2} = k/2.
\]
\end{enumerate}
प्रत्येक प्रकरणात विगमन गृहीतकाच्या आधारे $t$ मध्ये $[$ ची संख्या $< k/2$ आहे हे दाखवले. प्रबल विगमनाने विधान सिद्ध होते.
\end{proof}

\begin{prob}
अतिनीटस पदांचा संच पुढील रीतीने परिभाषित करा:
\begin{enumerate}
\item $\mathrm{a}$, $\mathrm{b}$, $\mathrm{c}$, $\mathrm{d}$ यांपैकी कोणतेही अक्षर अतिनीटस पद आहे.
\item $s$ हे अतिनीटस पद असेल, तर $[s]$ देखील अतिनीटस पद आहे.
\item $s_1$ आणि $s_2$ अतिनीटस पदे असतील, तर $[s_1 \circ s_2]$ देखील अतिनीटस पद आहे.
\item इतर काहीही अतिनीटस पद नाही.
\end{enumerate}
~$n$ लांबीच्या अतिनीटस पदात~$t$ मध्ये $[$ ची संख्या $\le n/2 +1$ आहे हे दाखवा.
\end{prob}

\end{document}
